\documentclass[10pt,a4paper]{amsart}
\usepackage{amsmath,amssymb,amsthm,graphicx}
% Retained minimal preamble for the Stacks Project chapter sites-cohomology.tex (Cohomology on Sites).
% Verbatim excerpt of the chapter preamble: Xy-pic setup (native preamble lines 14-18) and the native
% theorem declarations with the mathematical macros (native preamble lines 154-217). Copied mechanically;
% nothing added, removed or redefined. The wrapper class and the amsmath/amssymb/amsthm/graphicx base
% packages are supplied by the assembler, so the class/externaldocument part of the native preamble is omitted.
% For commutative diagrams we use Xy-pic
\usepackage[all]{xy}

% We use 2cell for 2-commutative diagrams.
\xyoption{2cell}
% Theorem environments.
%
\theoremstyle{plain}
\newtheorem{theorem}[subsection]{Theorem}
\newtheorem{proposition}[subsection]{Proposition}
\newtheorem{lemma}[subsection]{Lemma}

\theoremstyle{definition}
\newtheorem{definition}[subsection]{Definition}
\newtheorem{example}[subsection]{Example}
\newtheorem{exercise}[subsection]{Exercise}
\newtheorem{situation}[subsection]{Situation}

\theoremstyle{remark}
\newtheorem{remark}[subsection]{Remark}
\newtheorem{remarks}[subsection]{Remarks}

\numberwithin{equation}{subsection}

% Macros
%
\def\lim{\mathop{\mathrm{lim}}\nolimits}
\def\colim{\mathop{\mathrm{colim}}\nolimits}
\def\Spec{\mathop{\mathrm{Spec}}}
\def\Hom{\mathop{\mathrm{Hom}}\nolimits}
\def\Ext{\mathop{\mathrm{Ext}}\nolimits}
\def\SheafHom{\mathop{\mathcal{H}\!\mathit{om}}\nolimits}
\def\SheafExt{\mathop{\mathcal{E}\!\mathit{xt}}\nolimits}
\def\Sch{\mathit{Sch}}
\def\Mor{\mathop{\mathrm{Mor}}\nolimits}
\def\Ob{\mathop{\mathrm{Ob}}\nolimits}
\def\Sh{\mathop{\mathit{Sh}}\nolimits}
\def\NL{\mathop{N\!L}\nolimits}
\def\CH{\mathop{\mathrm{CH}}\nolimits}
\def\proetale{{pro\text{-}\acute{e}tale}}
\def\etale{{\acute{e}tale}}
\def\QCoh{\mathit{QCoh}}
\def\Ker{\mathop{\mathrm{Ker}}}
\def\Im{\mathop{\mathrm{Im}}}
\def\Coker{\mathop{\mathrm{Coker}}}
\def\Coim{\mathop{\mathrm{Coim}}}

% Boxtimes
%
\DeclareMathSymbol{\boxtimes}{\mathbin}{AMSa}{"02}

%
% Macros for moduli stacks/spaces
%
\def\QCohstack{\mathcal{QC}\!\mathit{oh}}
\def\Cohstack{\mathcal{C}\!\mathit{oh}}
\def\Spacesstack{\mathcal{S}\!\mathit{paces}}
\def\Quotfunctor{\mathrm{Quot}}
\def\Hilbfunctor{\mathrm{Hilb}}
\def\Curvesstack{\mathcal{C}\!\mathit{urves}}
\def\Polarizedstack{\mathcal{P}\!\mathit{olarized}}
\def\Complexesstack{\mathcal{C}\!\mathit{omplexes}}
% \Pic is the operator that assigns to X its picard group, usage \Pic(X)
% \Picardstack_{X/B} denotes the Picard stack of X over B
% \Picardfunctor_{X/B} denotes the Picard functor of X over B
\def\Pic{\mathop{\mathrm{Pic}}\nolimits}
\def\Picardstack{\mathcal{P}\!\mathit{ic}}
\def\Picardfunctor{\mathrm{Pic}}
\def\Deformationcategory{\mathcal{D}\!\mathit{ef}}
\usepackage{cleveref}
\crefname{originallemma}{original lemma}{original lemmas}
\Crefname{originallemma}{Original lemma}{Original lemmas}
\makeatletter
\newlabel{sites-definition-quasi-compact-topos}{{0D07}{source}{}{sourcepointer-1}{}}
\newlabel{sites-definition-quasi-compact-topos@cref}{{[originallemma][0][]0D07}{[1][0][]source}{}{}{}}
\makeatother
\crefname{originallemma}{original lemma}{original lemmas}
\Crefname{originallemma}{Original lemma}{Original lemmas}
\makeatletter
\newlabel{sites-lemma-directed-colimits-global-sections}{{0GMR}{source}{}{sourcepointer-2}{}}
\newlabel{sites-lemma-directed-colimits-global-sections@cref}{{[originallemma][0][]0GMR}{[1][0][]source}{}{}{}}
\makeatother
\crefname{originallemma}{original lemma}{original lemmas}
\Crefname{originallemma}{Original lemma}{Original lemmas}
\makeatletter
\newlabel{sites-lemma-image-of-quasi-compact}{{0GMP}{source}{}{sourcepointer-3}{}}
\newlabel{sites-lemma-image-of-quasi-compact@cref}{{[originallemma][0][]0GMP}{[1][0][]source}{}{}{}}
\makeatother
\crefname{originallemma}{original lemma}{original lemmas}
\Crefname{originallemma}{Original lemma}{Original lemmas}
\makeatletter
\newlabel{lemma-colim-works-over-collection}{{0739}{source}{}{sourcepointer-4}{}}
\newlabel{lemma-colim-works-over-collection@cref}{{[originallemma][0][]0739}{[1][0][]source}{}{}{}}
\makeatother
\crefname{originallemma}{original lemma}{original lemmas}
\Crefname{originallemma}{Original lemma}{Original lemmas}
\makeatletter
\newlabel{lemma-cech-to-cohomology-sheaf-sets}{{079Z}{source}{}{sourcepointer-5}{}}
\newlabel{lemma-cech-to-cohomology-sheaf-sets@cref}{{[originallemma][0][]079Z}{[1][0][]source}{}{}{}}
\makeatother
\crefname{originallemma}{original lemma}{original lemmas}
\Crefname{originallemma}{Original lemma}{Original lemmas}
\makeatletter
\newlabel{lemma-cohomology-on-sheaf-sets}{{07A0}{source}{}{sourcepointer-6}{}}
\newlabel{lemma-cohomology-on-sheaf-sets@cref}{{[originallemma][0][]07A0}{[1][0][]source}{}{}{}}
\makeatother
\crefname{originallemma}{original lemma}{original lemmas}
\Crefname{originallemma}{Original lemma}{Original lemmas}
\makeatletter
\newlabel{lemma-cohomology-of-open}{{03F3}{source}{}{sourcepointer-7}{}}
\newlabel{lemma-cohomology-of-open@cref}{{[originallemma][0][]03F3}{[1][0][]source}{}{}{}}
\makeatother
\crefname{originallemma}{original lemma}{original lemmas}
\Crefname{originallemma}{Original lemma}{Original lemmas}
\makeatletter
\newlabel{sites-lemma-localize-topos-site}{{0791}{source}{}{sourcepointer-8}{}}
\newlabel{sites-lemma-localize-topos-site@cref}{{[originallemma][0][]0791}{[1][0][]source}{}{}{}}
\makeatother
\title{Filtered colimits of abelian sheaves commute with cohomology under a generating set of quasi-compact sheaves}
\author{The Stacks Project Authors}
\date{stacks-project:sites-cohomology.tex:local-snapshot-20261011; GNU FDL 1.2 or later}
\begin{document}
\maketitle
\noindent\textit{Editorial note.} All delivered mathematical text is the author's own, copied from the pinned Stacks Project chapter file sites-cohomology.tex; no statement or proof step was rewritten, reordered or removed.\par\medskip
\noindent\textit{Editorial note.} The proof is the author's complete argument: the base case, the reduction to the vanishing of a single degree, the spectral sequence computation and the verification that the localized topoi satisfy the hypotheses are all inside the counted span.\par\medskip
\noindent\textit{Editorial note.} Every labelled reference occurring inside the counted spans points to a chapter of the Project that is not part of this card's source tree. Each is carried in assembly.external\_references as a non-load-bearing navigation pointer whose source location is the invocation inside the counted span, so no load-bearing mathematics is replaced by a stub.\par\medskip
\noindent\textit{Editorial note.} The checked chapter list for same-name collisions is the set of stacks chapters staged in this repository: chow, cohomology, divisors, etale-cohomology, flat, more-algebra, more-morphisms, morphisms, spaces-more-morphisms, varieties together with this card's chapter sites-cohomology.\par\medskip
\noindent\textit{Editorial note.} The preamble is the chapter's own retained minimal preamble (sources/stacks-sites-cohomology/retained-preamble.tex), a verbatim excerpt of the native preamble: the Xy-pic setup and the native theorem declarations with the mathematical macros.\par\medskip
\noindent\textit{Allowed standard prerequisites (editorial).}\begin{itemize}
\item Filtered colimits of abelian sheaves and their exactness. Same chapter, the description of filtered colimits of abelian sheaves and the exactness of filtered colimits of abelian groups. Supplies the diagrams to which the statement applies and the exactness used at the base of the induction.
\item Quasi-compact object of a topos and the quasi-compactness hypothesis (5). Sites, the definition of a quasi-compact object of a topos, recalled in the counted statement. Supplies the finiteness that makes the base case of the induction work.
\item Cohomology of a topos computed from a set of sheaves and the spectral sequence E\_1\textasciicircum{}\{p,q\} = H\textasciicircum{}q(K\textasciicircum{}\{p+1\}, I). Same chapter, the Cech-to-cohomology spectral sequence attached to a surjection K -> * and the identification H\textasciicircum{}q(K, F) = H\textasciicircum{}q(C/K, j\textasciicircum{}\{-1\}F). Supplies the spectral sequence with which the induction step reduces to the vanishing of H\textasciicircum{}\{p0+1\} on a filtered colimit of injective sheaves.
\item Localization of a topos at an object and the induction hypothesis in the localized topos. Same chapter, the topos C/K and the fact that the hypotheses (1)-(5) are inherited by it. Supplies the invocation of the induction hypothesis in the last paragraph of the proof.
\end{itemize}\medskip
\ifdefined\hypertarget\hypertarget{sourcepointer-1}{}\fi
\noindent\textit{Source reference.} Original lemma 0D07: sources/stacks-sites-cohomology/pinned-original-source.tex, source lines 2643--2643. Navigation pointer invoked at proof line 2643; the counted span contains the pointer only and the cited result is not reproduced inside it, so the pointer is registered as a non-load-bearing navigation stub in the established form for Stacks material. No bibliography citation is involved.\par\medskip
\ifdefined\hypertarget\hypertarget{sourcepointer-2}{}\fi
\noindent\textit{Source reference.} Original lemma 0GMR: sources/stacks-sites-cohomology/pinned-original-source.tex, source lines 2658--2658. Navigation pointer invoked at proof line 2658; the counted span contains the pointer only and the cited result is not reproduced inside it, so the pointer is registered as a non-load-bearing navigation stub in the established form for Stacks material. No bibliography citation is involved.\par\medskip
\ifdefined\hypertarget\hypertarget{sourcepointer-3}{}\fi
\noindent\textit{Source reference.} Original lemma 0GMP: sources/stacks-sites-cohomology/pinned-original-source.tex, source lines 2669--2669. Navigation pointer invoked at proof line 2669; the counted span contains the pointer only and the cited result is not reproduced inside it, so the pointer is registered as a non-load-bearing navigation stub in the established form for Stacks material. No bibliography citation is involved.\par\medskip
\ifdefined\hypertarget\hypertarget{sourcepointer-4}{}\fi
\noindent\textit{Source reference.} Original lemma 0739: sources/stacks-sites-cohomology/pinned-original-source.tex, source lines 2518--2518. Navigation pointer invoked at proof line 2682; the counted span contains the pointer only and the cited result is not reproduced inside it, so the pointer is registered as a non-load-bearing navigation stub in the established form for Stacks material. No bibliography citation is involved.\par\medskip
\ifdefined\hypertarget\hypertarget{sourcepointer-5}{}\fi
\noindent\textit{Source reference.} Original lemma 079Z: sources/stacks-sites-cohomology/pinned-original-source.tex, source lines 2073--2073. Navigation pointer invoked at proof line 2694; the counted span contains the pointer only and the cited result is not reproduced inside it, so the pointer is registered as a non-load-bearing navigation stub in the established form for Stacks material. No bibliography citation is involved.\par\medskip
\ifdefined\hypertarget\hypertarget{sourcepointer-6}{}\fi
\noindent\textit{Source reference.} Original lemma 07A0: sources/stacks-sites-cohomology/pinned-original-source.tex, source lines 2097--2097. Navigation pointer invoked at proof line 2697; the counted span contains the pointer only and the cited result is not reproduced inside it, so the pointer is registered as a non-load-bearing navigation stub in the established form for Stacks material. No bibliography citation is involved.\par\medskip
\ifdefined\hypertarget\hypertarget{sourcepointer-7}{}\fi
\noindent\textit{Source reference.} Original lemma 03F3: sources/stacks-sites-cohomology/pinned-original-source.tex, source lines 608--608. Navigation pointer invoked at proof line 2701; the counted span contains the pointer only and the cited result is not reproduced inside it, so the pointer is registered as a non-load-bearing navigation stub in the established form for Stacks material. No bibliography citation is involved.\par\medskip
\ifdefined\hypertarget\hypertarget{sourcepointer-8}{}\fi
\noindent\textit{Source reference.} Original lemma 0791: sources/stacks-sites-cohomology/pinned-original-source.tex, source lines 2731--2731. Navigation pointer invoked at proof line 2731; the counted span contains the pointer only and the cited result is not reproduced inside it, so the pointer is registered as a non-load-bearing navigation stub in the established form for Stacks material. No bibliography citation is involved.\par\medskip
\begin{lemma}
\label{lemma-colim-global}
Let $\mathcal{C}$ be a site. Let $S \subset \Ob(\Sh(\mathcal{C}))$
be a subset. Denote $*$ the final object of $\Sh(\mathcal{C})$. Assume
\begin{enumerate}
\item for some $K \in S$ the map $K \to *$ is surjective,
\item given a surjective map of sheaves $\mathcal{F} \to K$ with $K \in S$
there exists a $K' \in S$ and a map $K' \to \mathcal{F}$ such
that the composition $K' \to K$ is surjective,
\item given $K, K' \in S$ there is a surjection $K'' \to K \times K'$
with $K'' \in S$,
\item given $a, b : K \to K'$ with $K, K' \in S$ there exists a
surjection $K'' \to \text{Equalizer}(a, b)$ with $K'' \in S$, and
\item every $K \in S$ is quasi-compact
(Sites, Definition \ref{sites-definition-quasi-compact-topos}).
\end{enumerate}
Then for all $p \geq 0$ the map
$$
\colim_\lambda H^p(\mathcal{C}, \mathcal{F}_\lambda)
\longrightarrow
H^p(\mathcal{C}, \colim_\lambda \mathcal{F}_\lambda)
$$
is an isomorphism for every filtered diagram
$\Lambda \to \textit{Ab}(\mathcal{C})$, $\lambda \mapsto \mathcal{F}_\lambda$.
\end{lemma}

\begin{proof}
We will prove this by induction on $p$. The base case $p = 0$
follows from
Sites, Lemma \ref{sites-lemma-directed-colimits-global-sections} part (4).
We check the assumptions hold, but we urge the reader to skip this part.
Suppose $\mathcal{F} \to *$ is surjective.
Choose $K \in S$ and $K \to *$ surjective as in (1). Then
$\mathcal{F} \times K \to K$ is surjective. Choose
$K' \to \mathcal{F} \times K$ with $K' \in S$ and $K' \to K$
surjective as in (2).
Then there is a map $K' \to \mathcal{F}$ and $K' \to *$ is surjective.
Hence 
Sites, Lemma \ref{sites-lemma-directed-colimits-global-sections}
assumption (4)(a) is satisfied.
By Sites, Lemma \ref{sites-lemma-image-of-quasi-compact}, assumptions
(3) and (5) we see that $K \times K$ is quasi-compact
for all $K \in S$.
Hence 
Sites, Lemma \ref{sites-lemma-directed-colimits-global-sections}
assumption (4)(b) is satisfied.
This finishes the proof of the base case.

\medskip\noindent
Induction step. Assume the result holds for
$H^p$ for $p \leq p_0$ and for all topoi $\Sh(\mathcal{C})$
such that a set $S \subset \Ob(\Sh(\mathcal{C}))$ can be found
satisfying (1) -- (5). Arguing exactly as in the proof of
Lemma \ref{lemma-colim-works-over-collection}
we see that it suffices to show: given a filtered colimit
$\mathcal{I} = \colim \mathcal{I}_\lambda$ with
$\mathcal{I}_\lambda$ injective abelian sheaves, we have
$H^{p_0 + 1}(\mathcal{C}, \mathcal{I}) = 0$.
Choose $K \to *$ surjective with $K \in S$ as in (1).
Denote $K^n$ the $n$-fold self product of $K$.
Consider the spectral sequence
$$
E_1^{p, q} = H^q(K^{p + 1}, \mathcal{I}) \Rightarrow
H^{p + q}(*, \mathcal{I}) = H^{p + q}(\mathcal{C}, \mathcal{I})
$$
of Lemma \ref{lemma-cech-to-cohomology-sheaf-sets}. Recall that
$H^q(K^{p + 1}, \mathcal{F}) = H^q(\mathcal{C}/K^{p + 1}, j^{-1}\mathcal{F})$,
for any abelian sheaf on $\mathcal{C}$, see
Lemma \ref{lemma-cohomology-on-sheaf-sets}. We have
$j^{-1}\mathcal{I} = \colim j^{-1}\mathcal{I}_\lambda$
as $j^{-1}$ commutes with colimits. The restrictions
$j^{-1}\mathcal{I}_\lambda$ are injective abelian sheaves
on $\mathcal{C}/K^{p + 1}$ by Lemma \ref{lemma-cohomology-of-open}.
Below we will show that the induction hypothesis applies
to $\mathcal{C}/K^{p + 1}$ and hence we see that
$H^q(K^{p + 1}, \mathcal{I}) = \colim H^q(K^{p + 1}, \mathcal{I}_\lambda) = 0$
for $q < p_0 + 1$ (vanishing as $\mathcal{I}_\lambda$ is injective).
It follows that
$$
H^{p_0 + 1}(\mathcal{C}, \mathcal{I}) =
H^{p_0 + 1}\left(\ldots \to
H^0(K^{p_0}, \mathcal{I}) \to
H^0(K^{p_0 + 1}, \mathcal{I}) \to
H^0(K^{p_0 + 2}, \mathcal{I}) \to \ldots\right)
$$
Again using the induction hypothesis, the complex depicted on the right
hand side is the colimit over $\Lambda$ of the complexes
$$
\ldots \to
H^0(K^{p_0}, \mathcal{I}_\lambda) \to
H^0(K^{p_0 + 1}, \mathcal{I}_\lambda) \to
H^0(K^{p_0 + 2}, \mathcal{I}_\lambda) \to \ldots
$$
These complexes are exact as $\mathcal{I}_\lambda$ is an
injective abelian sheaf (follows from the spectral sequence
for example). Since filtered colimits are exact in the category
of abelian groups we obtain the desired vanishing.

\medskip\noindent
We still have to show that the induction hypothesis applies
to the site $\mathcal{C}/K^n$ for all $n \geq 1$. Recall
that $\Sh(\mathcal{C}/K^n) = \Sh(\mathcal{C})/K^n$, see
Sites, Lemma \ref{sites-lemma-localize-topos-site}.
Thus we may work in $\Sh(\mathcal{C})/K^n$.
Denote $S_n \subset \Ob(\Sh(\mathcal{C}/K^n)$ the set
of objects of the form $K' \to K^n$. We check each property in turn:
\begin{enumerate}
\item By (3) and induction there exists a surjection
$K' \to K^n$ with $K' \in S$. Then $(K' \to K^n) \to (K^n \to K^n)$
is a surjection in $\Sh(\mathcal{C})/K^n$ and $K^n \to K^n$
is the final object of $\Sh(\mathcal{C})/K^n$. Hence (1) holds for $S_n$,
\item Property (2) for $S_n$ is an immediate consequence
of (2) for $S$.
\item Let $a : K_1 \to K^n$ and $b : K_2 \to K^n$ be in $S_n$.
Then $(K_1 \to K^n) \times (K_2 \to K^n)$ is the object
$K_1 \times_{K^n} K_2 \to K^n$ of $\Sh(\mathcal{C})/K^n$.
The subsheaf $K_1 \times_{K^n} K_2 \subset K_1 \times K_2$ is the
equalizer of $a \circ \text{pr}_1$ and $b \circ \text{pr}_2$.
Write $a = (a_1, \ldots, a_n)$ and $b = (b_1, \ldots, b_n)$.
Pick $K_3 \to K_1 \times K_2$ surjective with $K_3 \in S$;
this is possibly by assumption (3) for $\mathcal{C}$.
Pick
$$
K_4
\longrightarrow
\text{Equalizer}(K_3 \to K_1 \times K_2 \xrightarrow{a_1, b_1} K)
$$
surjective with $K_4 \in S$.
This is possible by assumption (4) for $\mathcal{C}$.
Pick
$$
K_5
\longrightarrow
\text{Equalizer}(K_4 \to K_1 \times K_2 \xrightarrow{a_2, b_2} K)
$$
surjective with $K_5 \in S$.
Again this is possible. Continue in this fashion until we get
$$
K_{3 + n}
\longrightarrow
\text{Equalizer}(K_{2 + n} \to K_1 \times K_2 \xrightarrow{a_n, b_n} K)
$$
surjective with $K_{3 + n} \in S$. By construction
$K_{3 + n} \to K_1 \times_{K^n} K_2$
is surjective. Hence $(K_{3 + n} \to K^n)$ is in $S_n$ and
surjects onto the product $(K_1 \to K^n) \times (K_2 \to K^n)$.
Thus (3) holds for $S_n$.
\item Property (4) for $S_n$ is an immediate consequence of
property (4) for $S$.
\item Property (5) for $S_n$ is a consequence of
property (5) for $S$. Namely, an object $\mathcal{F} \to K^n$ of
$\Sh(\mathcal{C})/K^n$ corresponds to a quasi-compact object of
$\Sh(\mathcal{C}/K^n)$ if and only if $\mathcal{F}$ is a
quasi-compact object of $\Sh(\mathcal{C})$.
\end{enumerate}
This finishes the proof of the lemma.
\end{proof}

\end{document}
